
\documentclass{../concepts/prepareprotokol} %všechny balíčky pro protokol + titulní strana
\usepackage{graphicx}
\usepackage{booktabs}
\usepackage{amsmath}
\usepackage{siunitx}
\sisetup{inter-unit-product =$\cdot$}
\sisetup{separate-uncertainty=true}

\praktikum{II} %číslo praktika (I, II nebo III)
\autor{David Němec} %jméno autora
\datum{27.\,10.\,2025} %datum měření
\cislo{28} %číslo úlohy
\nazev{Charakteristické vlastnosti solárních článků} %název úlohy

\begin{document}
\maketitle

% ========================== PRACOVNÍ ÚKOL =========================
\section{Pracovní úkol}
\begin{enumerate}
    \item S použitím solárního simulátoru změřte na vybraném fotovoltaickém článku (FV) napětí naprázdno $U_{OC}$ a proud ve zkratu $I_{SC}$ v závislosti na intenzitě dopadajícího záření. Zpracujte graficky přímo v praktiku závislosti $U_{OC}$ a $I_{SC}$ na intenzitě dopadajícího záření.

    \item Pro tři vybrané intenzity záření slunečního simulátoru 0.25, 0.5 a 1.0 Sun určete z naměřených hodnot $U_{OC}$ a $I_{SC}$ optimální odpor zátěže $R_{L,opt}$ a z něho rozsah odporu zátěže $R_L$ pro měření zátěžových charakteristik. Rozsah $R_L$ volte tak, aby minimální odpor $R_L$ byl o řád menší než odpor v bodě maximálního výkonu a maximální odpor $R_L$ byl o řád větší než odpor v bodě maximálního výkonu.

    \item Pro tři vybrané intenzity záření slunečního simulátoru 0,25, 0,50 a 1,0 Sun změřte zátěžové charakteristiky fotovoltaického článku v rozsahu odporu zátěže $R_L$ určeném v bodě 2. Graficky zpracujte závislost proudu I a výkonu P na napětí U fotovoltaického článku.

    \item Z naměřených závislostí určete základní parametry FV článku, napětí $U_{mpp}$, proud $I_{mpp}$ a odpor $R_{mpp}$ článku v bodě maximálního výkonu, maximální výkon článku, faktor plnění $FF$ a účinnost $\eta$. Srovnejte a diskutujte účinnost fotovoltaického článku pro různou intenzitu osvětlení. Srovnejte optimální odpor zátěže $R_{L,opt}$ a zátěž fotovoltaického článku v bodě maximálního výkonu $R_{mpp}$.

    \item S použitím adaptéru s instalovanou LED diodou změřte závislost napětí naprázdno $U_{OC}$ a proudu ve zkratu $I_{SC}$ na intenzitě osvětlení pro nízké hodnoty osvětlení FV článku.

    \item Z naměřených závislostí určete následující parametry fotovoltaického článku, sériový odpor $R_s$, paralelní odpor $R_p$, saturační proud $I_0$ a faktor ideality $n$.
\end{enumerate}

% ============================= TEORIE ===============================
\section{Teorie}
\subsection{Kapacita deskového kondenzátoru}


\subsection{Metoda měření}

Zkoumaný fotovoltaický článek byl osvětlován solárním simulátorem s nastavitelnou intenzitou záření od 0.10 po 1.00 Sun s přesností na jednu setinu nominální intenzity slunečního záření. Pro odstínění okolního osvětlení byl použit černý válcový kryt. Protože se fotovoltaický článek dopadajícím zářením zahříval, bylo nutné jej chladit Peltiérovým článkem. Chlazení řídil termostat, který při překročení mezní teploty (cca 21.3 °C) sepnul obvod s Peltiérovým článkem, čímž umožnil chlazení fotovoltaického článku. Při podchlazení pod cílovou teplotu 20.5~°C termostat obvod rozpojil a chlazení článku bylo přerušeno. Schéma chladícího obvodu je na obrázku \ref{fig:schema}.

Během měření 1. úkolu byl obvod zapojen podle obrázku \ref{fig:schema1}: fotovoltaický článek byl napojen na přepínač, jímž se ke článku střídavě připojoval ampérmetr (pro měření zkratového proudu $I_{SC}$) nebo voltmetr (pro měření napětí otevřeného obvodu $U_{OC}$. Zkratový proud i napětí otevřeného obvodu byly měřeny pro deset intenzit záření od 0.1 Sun po 1.0 Sun, vždy bylo vyčkáno, než se hodnota trochu ustálí.

Zátěžové charakteristiky článku byly prováděny pomocí odporové dekády připojené paralelně k voltmetru podle schématu \ref{fig:schema2}. Měřeno bylo napětí na odporové dekádě, proud byl vypočten z Ohmova zákona (\ref{eq:ohm}).

Následně byl proměřován zkratový proud a napětí otevřeného obvodu pro nízké hodnoty osvětlení. Obvod byl zapojen opět podle schématu \ref{fig:schema1}. Článek byl osvětlován diodou umístěnou na spodní straně krytu, který se dal shora nasadit na válcový kryt kolem fotovoltaického článku.

\begin{figure}[htbp]
    \centering
    \includegraphics[width=0.5\linewidth]{uloha28/schema.png}
    \caption{Schéma zapojení Peltiérova článku (TEC) [2]}
    \label{fig:schema}
\end{figure}

\begin{figure}[htbp]
    \centering
    \includegraphics[width=0.5\linewidth]{uloha28/schema1.png}
    \caption{Schéma zapojení pro měření zkratového proudu $I_{SC}$ a napětí otevřeného obvodu $U_{OC}$ fotovoltaického článku s přepínačem [1]}
    \label{fig:schema1}
\end{figure}

\begin{figure}[htbp]
    \centering
    \includegraphics[width=0.5\linewidth]{uloha28/schema1.png}
    \caption{Schéma zapojení pro měření VA charakteristiky fotovoltaického článku se zátěží $R_L$ se známou hodnotou odporu a multimetrem [1]}
    \label{fig:schema2}
\end{figure}

\begin{figure}[htbp]
    \centering
    \includegraphics[width=0.5\linewidth]{uloha28/schema1.png}
    \caption{Ekvivalentní obvod reálného solárního článku obsahující generátor fotoproudu a
diodu pro oddělení náboje, sériový odpor $R_S$ a paralelní odpor $R_P$ připojené k zátěži $R_L$ [1]}
    \label{fig:schema3}
\end{figure}

% ========================= MĚŘICÍ PŘÍSTROJE ==========================
\subsection{Měřicí přístroje a jejich nejistoty}
\begin{enumerate}
    \item Multimetr Fluke 8808:
    
     Multimetrem byl měřen zkratový proud a napětí otevřeného obvodu v úkolech 1 a 5, a napětí na odporové dekádě v úkolu 3. Přesnost měření proudu je z manuálu k použití pro rozsah 200 $\mu$A 0.03\,\% hodnoty + 0.005\,\% rozsahu, pro rozsah 20 mA 0.04\,\% hodnoty + 0.02\,\% rozsahu a pro rozsah 200 mA 0.03\,\% hodnoty + 0.008\,\% rozsahu. Přesnost měření napětí je pro rozsah 200 mV 0.015\,\% hodnoty + 0.004\,\% rozsahu a pro rozsah 2 V 0.015\,\% hodnoty + 0.003\,\% rozsahu.

    \item Odporová dekáda:

     Odporová dekáda byla použita jako zátěžový odpor pro fotovoltaický článek, přesnost nastavené hodnoty odporu je ze specifikací přístroje 1\,\%.

\end{enumerate}
%\def\figurename{Obr\’azek}%
\def\figurename{Graf}%
\setcounter{figure}{0}

% ======================== VÝSLEDKY MĚŘENÍ ===========================
\section{Výsledky měření}


\begin{figure}
    \centering
    \includegraphics[width=0.5\linewidth]{uloha27/GRAF1.png}
    \caption{Závislost zkratového proudu $I_{SC}$ na napětí otevřeného obvodu $U_{OC}$. Vlevo se stejnými měřítky v obou osách, vpravo s logaritmickým měřítkem osy y.}
    \label{fig:desetintenzit}
\end{figure}


% ============================ DISKUZE ===============================
\section{Diskuze}


% ============================= ZÁVĚR ================================
\section{Závěr}


% =========================== LITERATURA =============================
\section*{Literatura}
\begin{itemize}
\item 

[1] Kolektiv ZFP KVOF MFF UK: Měření dielektrických vlastností materiálů [online]. [cit. 26.10.2025], dostupné z https://physics.mff.cuni.cz/vyuka/zfp/\_media/zadani/texty/txt\_227.pdf

\item 

[2] Kolektiv ZFP KVOF MFF UK: Měření odporu [online]. [cit. 26.10.2025], dostupné \\
z https://physics.mff.cuni.cz/vyuka/zfp/\_media/zadani/texty/txt\_202.pdf

\item

[3] Kolektiv ZFP KVOF MFF UK: Měření dielektrických vlastností materiálů - pokyny k měření [online]. [cit. 26.10.2025], dostupné z https://physics.mff.cuni.cz/vyuka/zfp/\_media/zadani/pokyny/mereni\_227.pdf

\item 

[4] Kolektiv ZFP KVOF MFF UK: Měření dielektrických vlastností materiálů - fotodokumentace [online]. [cit. 26.10.2025], dostupné z https://physics.mff.cuni.cz/vyuka/zfp/\_media/zadani/foto/foto\_227.pdf

\item 

[5] J. Englich: Úvod do praktické fyziky I: Zpracování výsledků měření. 1. vyd. Praha: Matfyzpress, 2006

\item 

[6] Eite Tiesinga, Peter J. Mohr, David B. Newell, and Barry N. Taylor: Values of fundamental physics constants: Vacuum electric permittivity [online]. [cit. 27.10.2025], dostupné z https://physics.nist.gov/cgi-bin/cuu/Value?ep0

\item 

[7] Dinel s.r.o.: tabulka relativních permitivit [online]. [cit. 27.10.2025], dostupné z https://www.dinel.cz/ \\ \_file/AMIfv95QsvMopq8GA848o69r\_yFAVYgklqLiAoLME0ORLQPdCAIkRAeJvDDR90OQh8kq6T4bVT6 \\ 
dASBZ27h\_lIuHiZ3j-GhH8T6DphAz1Ui4CwS-441Bupe\_IVYYHw\_XkUqAHRt6U93WEiawC8m6jG7tzF1pt \\
cn0sQ/relativni-permitivity-cz\_n1.pdf

\item 

[8] Mikulčák a kolektiv: Matematické, fyzikální a chemické tabulky pro střední školy: relativní permitivita, Praha: Státní pedagogické nakladatelství, n. p., 1988

\end{itemize}
\end{document}